% ============================================================
\section{The Persistent Geometric State}
\label{sec:persistent-geometric-state}
% ============================================================

The persistence framework begins by distinguishing the direct gravitational
response associated with the contemporary stress--energy distribution from a
surviving hereditary geometric state generated during the earlier evolution of
that distribution. The central hypothesis is not merely that these two
contributions coexist. It is that newly generated curvature may be realized
within a spacetime geometry that already contains a surviving state produced by
prior baryonic evolution.

In the standard initial-value formulation of General Relativity, the metric and
its admissible initial and boundary data describe the complete gravitational
state. Earlier source evolution may influence the present geometry through
retarded propagation and curved-spacetime Green-function tails. Such hereditary
contributions are ordinarily expected to disperse, decay, or become incorporated
into a metric configuration that subsequently relaxes as the source evolves.

Within the persistence framework, this relaxation is not assumed to be complete
on every physical scale. A portion of the causal hereditary response may survive
for a finite interval after the baryonic configuration that generated it has
moved, reorganized, or otherwise changed. We refer to this surviving
contribution as the \emph{persistent geometric state}.

The persistent geometric state is not identified with an additional matter
distribution. It is instead treated as an effective dynamical contribution to
spacetime geometry whose present configuration depends upon prior baryonic
evolution, its own admissible initial data, and its subsequent causal
propagation and relaxation.

Its existence would imply that the gravitational environment encountered by
matter is influenced not only by the direct response associated with the
contemporary stress--energy tensor, but also by an inherited geometric state and,
in general, by the interaction between that inherited state and the curvature
being generated by the contemporary source.

% ------------------------------------------------------------
\subsection{Defining weak-field decomposition}
\label{subsec:persistent-state-decomposition}
% ------------------------------------------------------------

In the weak-field regime, define the ordinary baryonic metric response by
\begin{equation}
    g_{\mu\nu}^{(b)}
    \equiv
    \eta_{\mu\nu}
    +
    h_{\mu\nu}^{(b)},
    \label{eq:baryonic-metric-background}
\end{equation}
where $\eta_{\mu\nu}$ is the Minkowski background and
$h_{\mu\nu}^{(b)}$ is the direct Einstein response associated with the
contemporary baryonic stress--energy tensor.

At lowest order, the observable metric may be decomposed as
\begin{equation}
    g_{\mu\nu}^{\rm eff}
    =
    \eta_{\mu\nu}
    +
    h_{\mu\nu}^{(b)}
    +
    H_{\mu\nu},
    \label{eq:persistent-metric-decomposition}
\end{equation}
where $H_{\mu\nu}$ is the effective metric perturbation representing the
persistent geometric state.

Equation~\eqref{eq:persistent-metric-decomposition} is a weak-field bookkeeping
decomposition. It does not represent two independent solutions of the ordinary
Einstein equations sourced by the same matter distribution, and it is not an
assumption that the direct and inherited sectors remain exactly additive at all
orders.

The field $h_{\mu\nu}^{(b)}$ denotes the direct baryonic response obtained from
the ordinary Einstein sector in the absence of persistence. The field
$H_{\mu\nu}$ denotes the surviving hereditary state associated with prior source
evolution, relative baryonic passage, or nontrivial persistent-sector initial
data.

Accordingly, $H_{\mu\nu}$ must not act as a second Poisson field sourced merely
by the continued presence of static matter. A permanently static source
comoving with the relevant geometric-state frame must not continuously
regenerate a duplicate of its ordinary gravitational field.

At this stage, $H_{\mu\nu}$ is introduced phenomenologically as the reduced
weak-field state variable. It is not assumed to be a fundamental unconstrained
rank-two field. Section~\ref{sec:covariant-parent-theory} constructs the
persistent sector from a restricted covariant field content consisting of a
dynamical timelike geometric-state frame together with scalar and vector
persistence degrees of freedom.

% ------------------------------------------------------------
\subsection{Interaction between inherited and contemporary curvature}
\label{subsec:persistent-state-interaction}
% ------------------------------------------------------------

The linear decomposition in
Eq.~\eqref{eq:persistent-metric-decomposition} is not taken to exhaust the
physical hypothesis. If a persistent geometric state is already present when
the contemporary baryonic configuration generates curvature, then the direct
response is realized in a geometry that is not identical to the fully relaxed
state. The observable gravitational response may therefore depend jointly on
the contemporary Einstein perturbation and the inherited state.

We represent this possibility schematically by a metric response functional
\begin{equation}
    g_{\mu\nu}^{\rm eff}
    =
    \mathcal{F}_{\mu\nu}
    \!\left[
        h_{\alpha\beta}^{(b)},
        H_{\alpha\beta}
    \right],
    \label{eq:persistent-effective-metric-functional}
\end{equation}
with the recovery condition
\begin{equation}
    \mathcal{F}_{\mu\nu}
    \!\left[
        h_{\alpha\beta}^{(b)},0
    \right]
    =
    \eta_{\mu\nu}
    +
    h_{\mu\nu}^{(b)}.
    \label{eq:persistent-functional-gr-recovery}
\end{equation}

For weak persistence and weak cross-coupling, the functional may be expanded as
\begin{equation}
    g_{\mu\nu}^{\rm eff}
    =
    \eta_{\mu\nu}
    +
    h_{\mu\nu}^{(b)}
    +
    H_{\mu\nu}
    +
    \mathcal{I}_{\mu\nu}
    \!\left[
        h^{(b)},H
    \right]
    +
    \mathcal{O}\!\left(H^2,h^{(b)}H^2,\ldots\right),
    \label{eq:persistent-metric-interaction-expansion}
\end{equation}
where $\mathcal{I}_{\mu\nu}$ denotes the leading cross-coupling between the
contemporary Einstein response and the inherited geometric state.

To isolate genuine interaction rather than a redefinition of either sector, we
require
\begin{equation}
    \mathcal{I}_{\mu\nu}[h^{(b)},0]=0,
    \qquad
    \mathcal{I}_{\mu\nu}[0,H]=0.
    \label{eq:persistent-interaction-vanishing}
\end{equation}
Pure self-interactions internal to the persistent sector, if present in the
nonlinear completion, are conceptually distinct and are represented by the
higher-order terms in
Eq.~\eqref{eq:persistent-metric-interaction-expansion}.

The interaction term is not specified phenomenologically on a galaxy-by-galaxy
basis. Its tensor form and couplings must follow from the same universal
covariant parent theory that governs the persistent state. In particular, it
must respect the gravitational constraint structure, the universal coupling of
matter to the effective metric, causal propagation, and continuous recovery of
General Relativity.

The physically observable departure from the direct baryonic metric is thus
\begin{equation}
    \Delta g_{\mu\nu}^{\rm her}
    \equiv
    g_{\mu\nu}^{\rm eff}
    -
    g_{\mu\nu}^{(b)}
    =
    H_{\mu\nu}
    +
    \mathcal{I}_{\mu\nu}[h^{(b)},H]
    +
    \cdots .
    \label{eq:observable-hereditary-metric-residual}
\end{equation}
Consequently, an observed dark-matter-like gravitational residual is not assumed
to measure $H_{\mu\nu}$ alone. It may instead measure the combined effect of the
surviving state and its coupling to the contemporary baryonic response.

This distinction yields two complementary empirical comparisons. Systems or
regions with similar present baryonic configurations but different source
histories test whether different inherited states produce different present
responses. Conversely, systems or regions with similar reconstructed inherited
states but different contemporary baryonic configurations test whether the
observable effect depends on the interaction between the inherited state and
the current source.

% ------------------------------------------------------------
\subsection{Distinction from standard hereditary effects}
\label{subsec:persistent-state-distinction}
% ------------------------------------------------------------

The persistent geometric state introduced here should be distinguished from
three related but physically different effects already present in General
Relativity: ordinary retardation, curvature-induced Green-function tails, and
gravitational-wave memory.

\paragraph{Ordinary retardation.}

Retardation reflects the finite propagation speed of gravitational
disturbances. A change in the source at $x'$ can influence an observation event
$x$ only when $x'$ lies within the causal past of $x$. In this sense,
retardation determines when information about a changing source can arrive. By
itself, however, retardation does not imply that a distinct local geometric
state remains after the corresponding disturbance has passed.

\paragraph{Curvature-induced Green-function tails.}

On a curved background, the retarded Green function of the linearized Einstein
equations generally has support not only on the null cone but also within it.
In Hadamard form,
\begin{equation}
    G^{\rm ret}_{\mu\nu}{}^{\alpha'\beta'}(x,x')
    =
    U_{\mu\nu}{}^{\alpha'\beta'}(x,x')\,
    \delta_{+}\!\left(\sigma\right)
    +
    V_{\mu\nu}{}^{\alpha'\beta'}(x,x')\,
    \Theta_{+}\!\left(-\sigma\right),
    \label{eq:gr-hadamard-distinction}
\end{equation}
where $\sigma(x,x')$ is Synge's world function. The $U$ term describes direct
propagation on the null cone, whereas the $V$ term represents curvature-induced
scattering within the null cone. The resulting metric at $x$ may therefore
depend on an extended portion of the causal source history.

This hereditary dependence is already part of standard General Relativity. It
remains, however, part of the ordinary Einstein-sector solution for the metric
and does not by itself introduce an independently evolving geometric state
variable possessing its own relaxation law, correlation length, state
occupancy, or interaction law with the contemporary source.

\paragraph{Gravitational-wave memory.}

Gravitational-wave memory refers to a lasting or asymptotic change in suitable
observables after the passage of radiative gravitational fields. Although
gravitational-wave memory is a genuine hereditary phenomenon, it is not
identified here with the persistent geometric state. In particular, the
state-defining timelike congruence introduced below is not a
gravitational-wave-memory degree of freedom, and the persistence variables do
not represent the standard displacement, spin, or related radiative memory
observables.

\paragraph{Persistent geometric state.}

The present framework introduces a different hypothesis. The additional state
variable
\begin{equation}
    H_{\mu\nu}
\end{equation}
is treated as an independently evolving geometric sector whose realized
configuration depends upon prior baryonic evolution and its own subsequent
dynamics. It possesses independent initial data and, in the effective
description developed below, characteristic propagation, relaxation, and
correlation scales.

The additional hypothesis further permits the inherited state to affect the
realization of later curvature through the universal interaction structure in
Eq.~\eqref{eq:persistent-metric-interaction-expansion}. Thus the framework is not
restricted to a passive residual field that simply adds to the contemporary
Einstein response.

Accordingly, the distinction may be summarized schematically as
\begin{equation}
\begin{aligned}
    \text{retardation}
    &:\quad
    \text{finite causal propagation},
    \\
    \text{GR curvature tail}
    &:\quad
    \text{Einstein-sector response with interior-cone support},
    \\
    \text{GW memory}
    &:\quad
    \text{lasting radiative observable change},
    \\
    \text{persistent state}
    &:\quad
    \text{independently evolving hereditary sector},
    \\
    \text{interaction}
    &:\quad
    \text{inherited--contemporary coupling}.
\end{aligned}
\label{eq:five-way-distinction}
\end{equation}

The novelty proposed here is therefore not the existence of hereditary
gravitational response itself. General Relativity already possesses hereditary
response through retarded propagation and curvature-induced tails. The
additional hypothesis is that a portion of the source-generated geometric
response can occupy a dynamical state that survives ordinary relaxation over a
finite interval and subsequently participates in later gravitational evolution.

This distinction also prevents double counting. The direct Einstein response to
the contemporary baryonic stress--energy distribution remains contained in
$h_{\mu\nu}^{(b)}$, while $H_{\mu\nu}$ represents only the additional inherited
state. The interaction term
$\mathcal{I}_{\mu\nu}[h^{(b)},H]$ is nonzero only when both sectors are present
and therefore cannot be interpreted as a second instantaneous gravitational
response sourced by the same static baryonic density.

% ------------------------------------------------------------
\subsection{Defining requirements}
\label{subsec:persistent-state-requirements}
% ------------------------------------------------------------

The persistent geometric state and its interaction with contemporary curvature
must satisfy five basic requirements.

\paragraph{1. Causal evolution.}

The persistent state must propagate causally. A change in the baryonic source
cannot alter $H_{\mu\nu}$ outside the characteristic cones of the physical
persistence modes.

Causality is therefore a condition on the principal part and characteristic
structure of the persistence equations, rather than an assumption that may be
imposed only after a solution has been obtained. Any interaction term affecting
the observable metric must inherit the same causal structure and must not
provide an acausal channel between the present source and the persistent state.

\paragraph{2. Independent initial-value description.}

The persistent state must possess admissible initial data of its own. Its value
at a spacetime event cannot be assumed to be a local algebraic function of the
contemporary stress--energy tensor alone.

The realized field may depend upon
\begin{equation}
    H_{\mu\nu}
    =
    H_{\mu\nu}
    \!\left[
        T_{\alpha\beta}^{(b)}(x'),
        H_{\alpha\beta}^{\rm initial}
    \right],
    \label{eq:persistent-history-functional-schematic}
\end{equation}
where the dependence on earlier source events $x'$ is restricted by causal
propagation.

Equation~\eqref{eq:persistent-history-functional-schematic} is schematic. The
explicit retarded functional is derived later from the Green functions of the
effective persistence operator. The contemporary metric response may then
depend jointly on this realized inherited state and on the present baryonic
source through the interaction functional
$\mathcal{I}_{\mu\nu}[h^{(b)},H]$.

\paragraph{3. Finite relaxation.}

Persistence does not imply permanent or unlimited accumulation. Earlier
contributions must be weighted by a causal response kernel and may decay,
disperse, dephase, or transfer energy into additional degrees of freedom.

The persistent state must therefore possess one or more finite response and
relaxation scales. For a representative mode, its free evolution should
approach the relaxed state rather than grow without bound.

The relevant distinction is thus
\begin{equation}
    \text{incomplete relaxation}
    \neq
    \text{absence of relaxation}.
    \label{eq:incomplete-not-absent-relaxation}
\end{equation}

Repeated source motion may generate a bounded persistent state, but it must not
lead automatically to secular divergence. The same requirement applies to
cross-coupling: interaction with contemporary curvature must not convert a
bounded inherited state into an uncontrolled instability.

\paragraph{4. Consistent interaction with contemporary curvature.}

The interaction between the inherited and contemporary sectors must be
universal and dynamically constrained. It may not be introduced as an arbitrary
system-dependent correction chosen to reproduce a desired rotation curve.

At minimum, the cross-coupling must satisfy
\begin{equation}
    \mathcal{I}_{\mu\nu}[h^{(b)},0]=0,
    \qquad
    \mathcal{I}_{\mu\nu}[0,H]=0,
\end{equation}
be compatible with the covariant constraint structure of the parent theory, and
preserve universal motion of matter in a single effective geometry.

The observable effect may therefore depend on both source history and the
contemporary baryonic configuration, but the law governing that dependence must
be the same for all physical systems.

\paragraph{5. Continuous recovery of General Relativity.}

The framework must recover the ordinary General Relativistic solution
continuously.

When the persistence coupling is removed, or when the persistent sector has
vanishing admissible initial data and receives no hereditary deposition, one
must have
\begin{equation}
    H_{\mu\nu}
    \longrightarrow
    0,
    \qquad
    \mathcal{I}_{\mu\nu}[h^{(b)},H]
    \longrightarrow
    0,
    \label{eq:persistent-state-gr-limit}
\end{equation}
so that
\begin{equation}
    g_{\mu\nu}^{\rm eff}
    \longrightarrow
    \eta_{\mu\nu}
    +
    h_{\mu\nu}^{(b)}.
    \label{eq:metric-gr-limit}
\end{equation}

The persistence framework therefore modifies neither the existence nor the
local form of the ordinary Einstein response in the persistence-free limit. It
introduces an additional effective hereditary sector and a possible universal
interaction with contemporary curvature, both of which decouple continuously
in the appropriate limit.

% ------------------------------------------------------------
\subsection{Physical interpretation}
\label{subsec:persistent-state-interpretation}
% ------------------------------------------------------------

The defining distinction may now be summarized schematically as
\begin{equation}
    \boxed{
    \begin{aligned}
    \text{present state}
    ={}&\ \text{direct Einstein response}
    \\
    &+\ \text{inherited geometric state}
    \\
    &+\ \text{inherited--contemporary interaction}.
    \end{aligned}
    }
    \label{eq:direct-persistent-interaction-response}
\end{equation}

Equation~\eqref{eq:direct-persistent-interaction-response} is not an independent
field equation. It summarizes the conceptual separation between the ordinary
baryonic response, the inherited state, and the leading nonlinear cross-coupling
between them.

At strictly linear order, the interaction term vanishes and the relation reduces
to
\begin{equation}
\begin{aligned}
    \text{present state}
    \simeq{}&\ \text{direct Einstein response}
    \\
    &+\ \text{surviving hereditary response}.
\end{aligned}
    \label{eq:linear-direct-plus-persistent-response}
\end{equation}
Thus the earlier additive picture is retained as the lowest-order limit rather
than discarded.

The persistent contribution need not be locally proportional to the direct
baryonic field. Depending upon its tensor structure, source history, propagation
history, and phase relative to the contemporary baryonic configuration, it may
reinforce, oppose, displace, or reshape the direct response. The cross-coupling
may further make the observable effect depend on how the contemporary baryonic
field is embedded in the inherited geometric state.

In the nonrelativistic weak-field limit, this means that the observed dynamical
acceleration may be represented schematically as
\begin{equation}
    \mathbf{a}_{\rm dyn}
    =
    \mathbf{a}_{b}
    +
    \mathbf{a}_{H}
    +
    \mathbf{a}_{\rm int}
    +
    \cdots,
    \label{eq:weak-field-acceleration-decomposition}
\end{equation}
so that the gravitational residual relative to the contemporary baryonic
prediction is
\begin{equation}
    \Delta\mathbf{a}
    \equiv
    \mathbf{a}_{\rm dyn}
    -
    \mathbf{a}_{b}
    =
    \mathbf{a}_{H}
    +
    \mathbf{a}_{\rm int}
    +
    \cdots .
    \label{eq:weak-field-hereditary-residual}
\end{equation}
The residual commonly attributed to a dark-matter contribution is therefore not
identified a priori with the persistent state alone. It is the observable
departure from the direct baryonic prediction and may contain both the inherited
state and its interaction with the present baryonic response.

Its observable effect is consequently not assumed to be equivalent to
\begin{itemize}
    \item a universal rescaling of Newton's constant;

    \item an additional static density profile;

    \item a second copy of the baryonic Newtonian potential;

    \item a local algebraic modification determined only by the present
    baryonic acceleration;

    \item or an arbitrary galaxy-specific correction to the gravitational law.
\end{itemize}

The universal content of the framework resides in the persistence evolution
law, the source couplings, the interaction law, and the propagation and
relaxation parameters. The realized persistent state of an individual physical
system depends upon its source history and admissible initial data, while the
observable gravitational effect depends on that realized state together with
the contemporary baryonic configuration.

Consequently, two systems with similar contemporary baryonic distributions need
not possess identical persistent geometric states if their histories contain
different patterns of orbital motion, gas accretion, star formation, radial
migration, mergers, stripping, secular reorganization, or environmental
interaction. Conversely, two systems with comparable inherited states need not
show identical residual accelerations if their contemporary baryonic
configurations differ and the cross-coupling is nonzero.

This yields a direct observational strategy. At fixed or closely matched present
baryonic configuration, differences in reconstructed source history test the
state dependence of the theory. At fixed or closely matched reconstructed
history, differences in present baryonic configuration test the interaction
term. A universal response law must account for both classes of comparison
without assigning a different persistence rule to each system.

This history dependence does not imply unrestricted freedom. The evolution
operator, deposition law, interaction structure, and couplings are universal.
Only the realized solution varies with the physical source history and present
source configuration, just as different initial and source data generate
different solutions of the same field equations.

% ------------------------------------------------------------
\subsection{Status of the persistent tensor}
\label{subsec:persistent-tensor-status}
% ------------------------------------------------------------

The tensor $H_{\mu\nu}$ should be interpreted carefully.

It is
\begin{itemize}
    \item the reduced weak-field state variable representing the inherited
    geometric contribution;

    \item an independently evolving hereditary state at the effective level;

    \item constrained by causal propagation, finite relaxation, and
    gravitational consistency;

    \item a contributor to, but not necessarily the entirety of, the observable
    gravitational residual;

    \item and ultimately constructed from a restricted covariant field content.
\end{itemize}

It is not
\begin{itemize}
    \item an additional material stress--energy tensor;

    \item a phenomenological dark-matter density inserted into the Einstein
    equations;

    \item a second instantaneous Poisson field sourced by the same static
    baryonic density;

    \item an unconstrained second metric introduced without a parent action;

    \item the interaction term itself;

    \item or the standard permanent gravitational-wave memory effect.
\end{itemize}

The distinction between state and observable residual is important. The field
$H_{\mu\nu}$ represents the inherited geometric state, whereas
\begin{equation}
    \Delta g_{\mu\nu}^{\rm her}
    =
    H_{\mu\nu}
    +
    \mathcal{I}_{\mu\nu}[h^{(b)},H]
    +
    \cdots
\end{equation}
represents the weak-field departure of the effective metric from the direct
baryonic metric. The latter is the quantity most closely connected to a
dark-matter-like observational residual.

The term \emph{persistent} refers to finite survival under incomplete
relaxation. It does not imply that $H_{\mu\nu}$ is permanent, nondynamical, or
independent of causal evolution. Nor does it imply that the inherited state is
observationally separable from the contemporary response without a dynamical
model of their coupling.

The purpose of the following section is to formulate the minimal postulates
governing this state and to construct the weakest causal transport equation
capable of representing incomplete geometric relaxation, together with a
consistent interaction with contemporary curvature, without duplicating the
ordinary direct Einstein field.
